\documentclass[11pt]{report}

\usepackage[pdftex]{geometry}
\geometry{a4paper,left=2.5cm,right=2.5cm,top=1.5cm,bottom=1.0cm,twoside}

\usepackage{amssymb}
\usepackage{enumerate}
\usepackage{graphicx}
	\graphicspath{./Comments/}
\usepackage[many]{tcolorbox}
\usepackage{natbib}
\usepackage{hyperref}
\usepackage{subfigure}
\usepackage{bm,bbm}
\usepackage{float}
\usepackage{siunitx}

\usepackage[utf8]{inputenc}
\usepackage[T1]{fontenc}
\usepackage{xcolor}
\newtcolorbox{reviewbox}[1]{
colback = white!5!white,
colframe = purple!75!black,
fonttitle=\bfseries,
title=#1
}

\newtcolorbox{responsebox}[1]{
	colback = white!5!white,
	colframe = white!75!black,
	fonttitle=\bfseries,
	title=#1
}
\newcommand{\assigned}[1]{\textcolor{blue}{\textbf{#1}}}
%%%%%%%%%%%%%%%%%%%%%%%%%%%%
\begin{document}


\begin{center}
\large{\textbf{Response Letter to Reviewer \# 2}}

\vglue 0.3cm

\huge{ Adaptive Model-Free Tsallis Entropy Estimation for Detecting Non-Fully-Developed Speckle\\ (TGRS-2025-10064)}
\end{center}

%\authors
\begin{center}
\textbf{Janeth Alpala,  Alejandro C.\ Frery, Paolo Gamba, Abraão D. C. Nascimento, Andrea A. Rey }
\end{center}

\date{\today}

%%%%%%%%%%%%%%%%%%%%%%%%%%%%

\vspace{0.5cm}
\noindent Dear Editor
\bigskip

\noindent We are grateful to the reviewers for the valuable comments and the time spent on checking our manuscript. 
We have addressed their comments towards improving the paper contents and presentation. 
Below we detail the changes and updates incorporated into the manuscript in this round of review.

\medskip



\vspace{0.5cm}

\begin{reviewbox}{Comment 1}

\textbf{Concerns Regarding Novelty and Incremental Contribution}

The manuscript is closely related to the authors' recent works, specifically the Shannon-entropy test ([18], [30]) and the Rényi-entropy test ([19]). In this submission, the main changes appear to be: (a) replacing Shannon/Rényi with Tsallis entropy and (b) adding an adaptive window selection step.

- Relationship to prior work

Moving from Shannon $\rightarrow$ Rényi $\rightarrow$ Tsallis is mathematically reasonable, but it may read as a straightforward extension unless the paper clearly demonstrates a distinct statistical or practical advantage of Tsallis entropy in this specific SAR setting. At the moment, the paper appears to follow a very similar testing pipeline to [18] and [19], with a different entropy definition.

- Recommendations

Please clearly state and quantitatively demonstrate why Tsallis entropy is preferable here. In particular, does the extra parameter $\lambda$ lead to a meaningful gain (e.g., better separation of texture levels, improved detection in certain heterogeneity regimes, or stronger robustness) compared with Rényi and Shannon approaches?

\end{reviewbox}

\begin{responsebox}{Response}

We thank the reviewer for this important observation. We revised the manuscript to clarify that the advantage of the proposed method is not merely the presence of an additional parameter, since the Rényi entropy also depends on an order parameter. Rather, the contribution lies in the way the Tsallis entropy uses the entropic index $\lambda$ to modify the sensitivity of the statistic to low-probability and tail events.

In SAR intensity data, the main differences between the fully developed speckle model and heterogeneous models appear in the tails of the distribution. For $\lambda<1$, the Tsallis entropy increases the influence of small-probability events, which helps emphasize departures from the Gamma model. This provides a practical statistical motivation for the selected value $\lambda=0.85$, which was also supported by the bias--MSE analysis reported in Section III-A.

\end{responsebox}

\begin{responsebox}{}



To demonstrate that this effect leads to a meaningful gain over Shannon and Rényi, we added a direct comparison in Section IV-E, entitled \emph{Comparative Benchmark Between Entropy-Based Tests}. The comparison uses the same Dublin SAR scene, ROIs, window settings, bootstrap configuration, and decision protocol for Shannon, Rényi, and Tsallis. The results in Table VI show that the bootstrap-corrected adaptive Tsallis test achieved the best overall performance, with the highest AUC, $F_1$, $\kappa$, and OA.

All three methods identify the main heterogeneous structures, but the Tsallis map provides a clearer separation of homogeneous water regions, with $p$-values closer to one, while preserving finer details in the urban and harbor areas.

\end{responsebox}

\vspace{1em}

\begin{reviewbox}{Comment 2}

\textbf{Missing benchmarking against entropy-based baselines}

In Section IV, the experiments mainly compare adaptive windows versus fixed windows. This is not sufficient to support the central claim about the Tsallis-based statistic.

- Missing baselines

Since Shannon- and Rényi-based tests already exist in [18] and [19], the paper should compare the proposed Tsallis test against those baselines directly.

- Isolating the source of improvement

Without a comparison such as ``Adaptive Tsallis vs.\ Adaptive Rényi vs.\ Adaptive Shannon'', it is not possible to tell whether improvements come from the Tsallis statistic or simply from adaptive windowing (which could be applied to any entropy measure).

- Recommendations

At minimum, please provide a threshold sweep (e.g., $p$-threshold from $0.01$ to $0.20$) for Tsallis, Rényi, and Shannon under the same windowing protocol and on the same ROIs (ideally for both fixed and adaptive windows). If space permits, ROC curves and/or precision--recall curves would provide a complementary, threshold-free view.

\end{reviewbox}

\begin{responsebox}{Response}

We agree with the reviewer. In the previous version, the experimental section mainly emphasized the effect of the adaptive windowing strategy within the Tsallis-based formulation. This did not fully isolate the contribution of the entropy functional itself.

To address this concern, we added a new entropy-level benchmark in Section IV-E. In this benchmark, Shannon, Rényi, and Tsallis were compared under the same protocol, using the Dublin SAR scene, the same ROIs, the same fixed $7\times7$ window, the same adaptive window range from $5\times5$ to $11\times11$, and the same decision rule based on $p_{ij}<0.05$. Configurations with and without bootstrap correction were also considered.

The revised manuscript now reports AUC, $F_1$, $\kappa$, and overall accuracy for all entropy-based tests. In addition, we added ROC curves as a threshold-free comparison and an $F_1$ threshold sweep using $p$-thresholds from $0.01$ to $0.20$.

The results show two consistent trends. First, the adaptive windowing strategy improves all three entropy-based tests. Second, within both fixed and adaptive settings, the Tsallis-based test provides the strongest overall performance. Therefore, the improvement cannot be attributed only to adaptive windowing; the Tsallis statistic itself contributes additional discriminatory power when combined with bootstrap correction and adaptive window selection.

These additions appear in the new subsection \emph{Comparative Benchmark Between Entropy-Based Tests}, together with Table~VI, Figure~15, Figure~16, and Figure~17.

\end{responsebox}


\vspace{1em}


\begin{reviewbox}{Comment 3}

\textbf{Calibration with adaptive window selection}

I am concerned about whether the reported $p$-values remain well calibrated when the window size is chosen adaptively (see the adaptive-windowing part around Eq.~(15)).

- Data-dependent window selection

The manuscript reports null behavior (Table I) based on Monte Carlo simulations for fixed sample sizes. In the proposed pipeline, however, the window size $W_{ij}$ is selected adaptively based on local homogeneity. Because the final test statistic is computed using a data-dependent window size, its null distribution (and therefore $p$-value calibration) may differ from the fixed-window setting; this point merits clarification and validation.

- Recommendations

Please discuss whether a fixed-window null distribution is an appropriate approximation for the adaptive-window statistic. Ideally, the calibration procedure should incorporate the adaptive selection mechanism (e.g., by simulating the full adaptive procedure under $H_0$ and estimating the null distribution of the final standardized statistic). If a full simulation is too expensive, a practical alternative is to calibrate separately for each admissible window size and to standardize per pixel using its final window size.

\end{reviewbox}

\begin{responsebox}{Response}

We agree that the behavior reported in Table I was obtained for fixed sample sizes, whereas in the adaptive procedure the final window size $W_{ij}$ is selected from the data. Therefore, the sampling behavior of the adaptive statistic does not necessarily coincide with that of a fixed-window statistic.

To address this point, we added a new subsection in Section III entitled \emph{Test statistic behavior under the null hypothesis and adaptive windows}. In this subsection, we discuss the effect of data-dependent window selection and evaluate the statistic under the complete adaptive procedure.

For this purpose, we generated a homogeneous $1000\times1000$ image under $\mathcal H_0$, with independent and identically distributed observations from $\Gamma_{\mathrm{SAR}}(10,16)$. We used $\lambda=0.85$, bootstrap size $B=100$, admissible window sizes $W\in\{5,7,9,11\}$, and $\eta=3$. To avoid edge effects, we analyzed an interior region obtained after removing a five-pixel guard region from each side. This guarantees that every pixel in the analyzed region can support the largest admissible window, namely $11\times11$, when required by the adaptive rule.

The experiment showed that, under $\mathcal H_0$, the adaptive procedure behaves very similarly to the fixed $W=11$ case in the interior region. This occurs because the adaptive procedure progressively expands the local window whenever the homogeneity criterion is satisfied, so the selected size tends to reach the maximum admissible value. We therefore added Figure~6 and Table~II, which respectively show the empirical densities and report the empirical mean and standard deviation of the statistic for each admissible window size.

For the computation of the adaptive $p$-value maps, we used the following standardization rule:
\[
z_{ij}=
\frac{S_{ij}-\mu_{\mathcal H_0}(W_{ij})}{\sigma_{\mathrm{obs}}},
\qquad
p_{ij}=2\Phi\bigl(-|z_{ij}|\bigr).
\]
This rule preserves the window-dependent centering through $\mu_{\mathcal H_0}(W_{ij})$, while using an observed-scale denominator better matched to the variability of the statistic field in SAR scenes.

We also applied the same standardization strategy in the comparative benchmark for the adaptive Shannon- and Rényi-based tests, so that all adaptive entropy-based methods were compared under the same protocol in Section IV-E.

\end{responsebox}

\end{document}

